%! Multiplying \& Dividing Radicals; Rationalizing — Unit 6, Lesson 6.3 — Group Activity
\documentclass[10pt]{article}
\usepackage{algebra2-article}
\usepackage{algebra2-boxes}
\newcommand{\TallMath}[1]{$\displaystyle #1\rule[-1.4em]{0pt}{3.2em}$}

\begin{document}

\pageheader{Unit 6, Lesson 6.3}{Group Activity}
\namepartnerperiod

\vspace{0.08in}

\begin{headlinebox}{forestbg}
\small\textbf{Making It Rational.} Irrational numbers are stubborn --- but multiply the right two
of them together and the radicals vanish completely. With your partner, hunt for the products that
come out \emph{rational}, and figure out what those problems have in common. Every final answer
must be in simplest radical form with \textbf{no radical in a denominator}. Assume all variables
represent \textbf{positive} real numbers.
\end{headlinebox}

\vspace{0.06in}

\begin{tcolorbox}[colback=white, colframe=black!40, title=\textbf{Tier R --- Remediate},
    fonttitle=\bfseries, arc=1.5mm, left=3mm, right=3mm, top=2mm, bottom=2mm, breakable]
\textbf{1. Multiply, then simplify.} Insides times insides, outsides times outsides. The first row
is done for you --- notice that the \emph{last} column is where the real work is.
\begin{center}
\renewcommand{\arraystretch}{1.7}
\begin{tabularx}{\linewidth}{@{}l X X@{}}
\hline
\textbf{Problem} & \textbf{One radical} & \textbf{Simplified} \\
\hline
$\sqrt{2}\cdot\sqrt{6}$      & $\sqrt{12}$ & $2\sqrt3$ \\
$\sqrt{2}\cdot\sqrt{8}$      & \blank{2.2cm} & \blank{2.2cm} \\
$\sqrt{3}\cdot\sqrt{6}$      & \blank{2.2cm} & \blank{2.2cm} \\
$\sqrt{5}\cdot\sqrt{5}$      & \blank{2.2cm} & \blank{2.2cm} \\
$2\sqrt{3}\cdot4\sqrt{2}$    & \blank{2.2cm} & \blank{2.2cm} \\
$\sqrt[3]{2}\cdot\sqrt[3]{4}$ & \blank{2.2cm} & \blank{2.2cm} \\
\hline
\end{tabularx}
\end{center}
Three rows above finished with \emph{no radical at all}. Which three, and what did the three
products have in common? \writeline

\vspace{3pt}
\textbf{2. Distribute.} Multiply the outside radical by \emph{each} term inside, then simplify each
piece.
\begin{center}
\renewcommand{\arraystretch}{2.0}
\begin{tabularx}{\linewidth}{@{}X X X@{}}
(a) \; $\sqrt3\left(\sqrt3+\sqrt5\right)=$ \blank{2.4cm} &
(b) \; $\sqrt2\left(\sqrt6-4\right)=$ \blank{2.4cm} &
(c) \; $\sqrt5\left(2\sqrt5+\sqrt{10}\right)=$ \blank{2.4cm} \\
\end{tabularx}
\end{center}

\vspace{3pt}
\textbf{3. What do I multiply by?} Fill in the form of $1$ you need, then finish.
\begin{center}
\renewcommand{\arraystretch}{2.2}
\begin{tabularx}{\linewidth}{@{}l c l c l@{}}
\hline
\textbf{Start} & & \textbf{Multiply by} & & \textbf{Answer} \\
\hline
$\dfrac{1}{\sqrt3}$ & $\times$ & $\dfrac{\sqrt3}{\sqrt3}$ & $=$ & $\dfrac{\sqrt3}{3}$ \\
$\dfrac{2}{\sqrt7}$ & $\times$ & \blank{1.4cm} & $=$ & \blank{2.0cm} \\
$\dfrac{5}{\sqrt5}$ & $\times$ & \blank{1.4cm} & $=$ & \blank{2.0cm} \\
$\dfrac{\sqrt3}{\sqrt2}$ & $\times$ & \blank{1.4cm} & $=$ & \blank{2.0cm} \\
\hline
\end{tabularx}
\end{center}
The third row simplifies further than you might expect --- check it.

\vspace{3pt}
\textbf{4. Conjugates.} Write each conjugate (same terms, opposite middle sign), then multiply the
pair.
\begin{center}
\renewcommand{\arraystretch}{1.7}
\begin{tabularx}{\linewidth}{@{}l X X@{}}
\hline
\textbf{Expression} & \textbf{Conjugate} & \textbf{Product} \\
\hline
$2+\sqrt{7}$          & \blank{2.6cm} & \blank{2.0cm} \\
$\sqrt{5}-3$          & \blank{2.6cm} & \blank{2.0cm} \\
$\sqrt{6}+\sqrt{2}$   & \blank{2.6cm} & \blank{2.0cm} \\
\hline
\end{tabularx}
\end{center}
Every product came out with no radical in it. Explain why that always happens. \writeline
\end{tcolorbox}

\begin{tcolorbox}[colback=white, colframe=black!40, title=\textbf{Tier A --- Approaching Proficiency},
    fonttitle=\bfseries, arc=1.5mm, left=3mm, right=3mm, top=2mm, bottom=2mm, breakable]
\textbf{1. Multiply and simplify completely.}
\begin{center}
\renewcommand{\arraystretch}{2.0}
\begin{tabularx}{\linewidth}{@{}X X@{}}
(a) \; $\sqrt6\cdot\sqrt{15}=$ \blank{2.4cm} &
(b) \; $3\sqrt2\cdot5\sqrt{10}=$ \blank{2.4cm} \\
(c) \; $\sqrt{2x}\cdot\sqrt{8x^{3}}=$ \blank{2.4cm} &
(d) \; $\sqrt[3]{9}\cdot\sqrt[3]{6}=$ \blank{2.4cm} \\
\end{tabularx}
\end{center}

\vspace{3pt}
\textbf{2. Expand and simplify.}
\begin{center}
\renewcommand{\arraystretch}{2.2}
\begin{tabularx}{\linewidth}{@{}X X@{}}
(a) \; $\sqrt3\left(\sqrt{12}-\sqrt6\right)=$ \blank{2.6cm} &
(b) \; $\left(4+\sqrt2\right)\left(3-\sqrt2\right)=$ \blank{2.6cm} \\
(c) \; $\left(\sqrt5+\sqrt3\right)^{2}=$ \blank{2.6cm} &
(d) \; $\left(2\sqrt7-3\right)\left(2\sqrt7+3\right)=$ \blank{2.6cm} \\
\end{tabularx}
\end{center}
Only one of (c) and (d) came out rational. Circle it, and say in a few words why the other did not.
\writeline

\vspace{3pt}
\textbf{3. Rationalize each denominator.}
\begin{center}
\renewcommand{\arraystretch}{2.4}
\begin{tabularx}{\linewidth}{@{}X X@{}}
(a) \; $\dfrac{8}{\sqrt6}=$ \blank{2.4cm} &
(b) \; $\dfrac{\sqrt5}{\sqrt{18}}=$ \blank{2.4cm} \\
(c) \; $\dfrac{3}{5+\sqrt2}=$ \blank{2.8cm} &
(d) \; $\dfrac{\sqrt3}{\sqrt7-\sqrt3}=$ \blank{2.8cm} \\
\end{tabularx}
\end{center}

\vspace{3pt}
\textbf{4. Find the mistake.} Each student made exactly one error. Name it and give the correct
answer.
\begin{center}
\renewcommand{\arraystretch}{1.9}
\begin{tabularx}{\linewidth}{@{}l X X@{}}
\hline
\textbf{Student work} & \textbf{What went wrong} & \textbf{Correct answer} \\
\hline
$\left(\sqrt5+\sqrt2\right)^{2}=5+2=7$ & \blank{4.0cm} & \blank{2.2cm} \\
$\dfrac{2}{\sqrt3}=\dfrac{2\sqrt3}{\sqrt3\cdot\sqrt3}=\dfrac{2\sqrt3}{\sqrt9}=\dfrac{2\sqrt3}{9}$
  & \blank{4.0cm} & \blank{2.2cm} \\
$\dfrac{5}{2+\sqrt3}=\dfrac{5\left(2+\sqrt3\right)}{\left(2+\sqrt3\right)\left(2+\sqrt3\right)}
  =\dfrac{10+5\sqrt3}{7+4\sqrt3}$ & \blank{4.0cm} & \blank{2.2cm} \\
\hline
\end{tabularx}
\end{center}

\vspace{3pt}
\textbf{5. A rectangle with irrational sides.} A rectangle measures $\left(3+\sqrt5\right)$ cm by
$\left(3-\sqrt5\right)$ cm.
\begin{enumerate}[label=(\alph*), leftmargin=1.9em, itemsep=4pt]
  \item Find its \textbf{area}: \blank{2.6cm} \quad Is that number rational or irrational?
        \blank{2.0cm}
  \item Find its \textbf{perimeter}: \blank{2.6cm} \quad Rational or irrational? \blank{2.0cm}
  \item Both sides are irrational, yet neither answer has a radical in it. Explain what happened
        in each case --- they are two \emph{different} reasons. \writelines{2}
  \item Now find the area of a \textbf{square} whose side is $\left(3+\sqrt5\right)$ cm:
        \blank{2.8cm} \\
        This one \emph{is} irrational. What is different about it? \writeline
\end{enumerate}
\end{tcolorbox}

\begin{tcolorbox}[colback=white, colframe=black!40, title=\textbf{Tier E --- Extension},
    fonttitle=\bfseries, arc=1.5mm, left=3mm, right=3mm, top=2mm, bottom=2mm, breakable]
\textbf{Part 1 --- Why the conjugate, and \emph{only} the conjugate.} You are going to rationalize
$\dfrac{1}{3-\sqrt5}$ three times, multiplying top and bottom by three different things. Only one
works.
\begin{enumerate}[label=(\alph*), leftmargin=1.9em, itemsep=5pt]
  \item Multiply by $\dfrac{\sqrt5}{\sqrt5}$. New denominator: \blank{2.6cm} \; Rational?
        \blank{1.2cm}
  \item Multiply by $\dfrac{3-\sqrt5}{3-\sqrt5}$. New denominator: \blank{2.6cm} \; Rational?
        \blank{1.2cm}
  \item Multiply by $\dfrac{3+\sqrt5}{3+\sqrt5}$. New denominator: \blank{2.6cm} \; Rational?
        \blank{1.2cm} \; Finish the problem: \blank{2.6cm}
  \item Now prove it in general. Expand $\left(a+\sqrt b\right)\left(a-\sqrt b\right)$ term by term
        and explain why the answer must be rational whenever $a$ and $b$ are rational.
        \writelines{2}
  \item \textbf{Lesson 3.4, revisited.} Back in the quadratics unit you multiplied $a+bi$ by
        $a-bi$ and got $a^{2}+b^{2}$. Today $\left(a+\sqrt b\right)\left(a-\sqrt b\right)=a^{2}-b$.
        Why is one a \emph{plus} and the other a \emph{minus}? \writeline
\end{enumerate}

\vspace{5pt}
\textbf{Part 2 --- Rationalizing the \emph{numerator} (this is a calculus problem).} Consider
\[
  E(x)=\frac{\sqrt{x+9}-3}{x}.
\]
\begin{enumerate}[label=(\alph*), leftmargin=1.9em, itemsep=5pt]
  \item Try to evaluate $E(0)$. What goes wrong? \blank{4.6cm}
  \item Multiply the top and bottom by $\sqrt{x+9}+3$ --- the conjugate of the \emph{numerator}.
        Show that the top becomes just $x$, and then that the whole expression reduces to
        \[
          E(x)=\frac{1}{\blank{3.0cm}}.
        \]
  \item Fill in the table using \emph{either} form (they are equal for every $x\ne0$). Round to
        four decimal places.
        \begin{center}
        \renewcommand{\arraystretch}{1.4}
        \begin{tabular}{|c|c|c|c|c|}
        \hline
        $x$ & $7$ & $1$ & $0.1$ & $0.01$ \\ \hline
        $E(x)$ & \blank{1.6cm} & \blank{1.6cm} & \blank{1.6cm} & \blank{1.6cm} \\ \hline
        \end{tabular}
        \end{center}
  \item The values are closing in on a number. Which one? \blank{1.8cm} \; Now put $x=0$ into
        \emph{your rationalized form} from (b): \blank{1.8cm}
  \item The original expression is undefined at $x=0$; the rewritten one is not, and it lands
        exactly on the number the table was approaching. Explain what rationalizing accomplished
        here. \writelines{2}
\end{enumerate}

\vspace{5pt}
\textbf{Part 3 --- Rationalizing when the index is not $2$.} Yesterday's rule --- \emph{the index
sets the group size} --- decides what you must multiply by.
\begin{enumerate}[label=(\alph*), leftmargin=1.9em, itemsep=5pt]
  \item $\dfrac{1}{\sqrt[3]{2}}$: you need \blank{1.0cm} copies of $2$ under the cube root, and you
        already have $1$, so multiply by $\dfrac{\blank{1.2cm}}{\blank{1.2cm}}$. Answer:
        \blank{2.0cm}
  \item $\dfrac{5}{\sqrt[3]{9}}$: \; multiply by \blank{1.4cm} \; Answer: \blank{2.2cm}
  \item $\dfrac{1}{\sqrt[4]{8}}$: \; multiply by \blank{1.4cm} \; Answer: \blank{2.2cm}
  \item State the rule in one sentence: to rationalize $\dfrac{1}{\sqrt[n]{a^{m}}}$, multiply by
        whatever \writeline
  \item \textbf{The same thing in exponent notation (Lesson 6.1).} Show that
        $\dfrac{1}{\sqrt[3]{2}}=2^{-1/3}$, and that $2^{-1/3}=\dfrac{2^{2/3}}{2}$. Does that agree
        with your answer to (a)? \writelines{2}
\end{enumerate}
\end{tcolorbox}

\end{document}
